\chapter{PENDAHULUAN}
\label{bab:pendahuluan}

\section{Latar Belakang}
\label{sec:latar_belakang}

Dummy latar belakang penelitian. Ganti dengan konten sebenarnya yang menjelaskan konteks, urgensi, dan motif penelitian.

Latar belakang harus menjawab pertanyaan: Mengapa topik ini penting? Apa masalah yang ingin dipecahkan? Bagaimana konteks ilmiah dan praktisnya?

\section{Rumusan Masalah}
\label{sec:rumusan_masalah}

Berdasarkan latar belakang yang telah diuraikan, rumusan masalah dalam penelitian ini adalah:

\begin{enumerate}
    \item Rumusan masalah pertama?
    \item Rumusan masalah kedua?
    \item Rumusan masalah ketiga?
\end{enumerate}

\section{Tujuan Penelitian}
\label{sec:tujuan_penelitian}

Sejalan dengan rumusan masalah di atas, tujuan penelitian ini adalah:

\begin{enumerate}
    \item Tujuan pertama yang ingin dicapai.
    \item Tujuan kedua yang ingin dicapai.
    \item Tujuan ketiga yang ingin dicapai.
\end{enumerate}

\section{Batasan Masalah}
\label{sec:batasan_masalah}

Agar penelitian tetap fokus dan memiliki metodologi yang kuat, penelitian ini menerapkan batasan masalah sebagai berikut:

\begin{enumerate}
    \item \textbf{Batasan pertama}: Deskripsi batasan pertama.
    \item \textbf{Batasan kedua}: Deskripsi batasan kedua.
    \item \textbf{Batasan ketiga}: Deskripsi batasan ketiga.
\end{enumerate}

\section{Manfaat Penelitian}
\label{sec:manfaat_penelitian}

Penelitian ini diharapkan dapat memberikan manfaat sebagai berikut:

\begin{enumerate}
    \item \textbf{Manfaat Teoretis}: Deskripsi kontribusi teoretis penelitian.
    \item \textbf{Manfaat Praktis}: Deskripsi kontribusi praktis penelitian.
\end{enumerate}

\section{Sistematika Penulisan}
\label{sec:sistematika_penulisan}

Penulisan tesis ini disusun dalam lima bab dengan sistematika sebagai berikut:

\begin{description}
    \item[BAB 1 -- PENDAHULUAN] Berisi latar belakang masalah, rumusan masalah, tujuan penelitian, batasan masalah, manfaat penelitian, dan sistematika penulisan.
    \item[BAB 2 -- TINJAUAN PUSTAKA] Menguraikan landasan teoretis yang mendukung penelitian.
    \item[BAB 3 -- METODOLOGI PENELITIAN] Menjelaskan rancangan eksperimen dan metodologi yang digunakan.
    \item[BAB 4 -- HASIL DAN PEMBAHASAN] Menyajikan hasil eksperimen dan analisis.
    \item[BAB 5 -- KESIMPULAN DAN SARAN] Berisi kesimpulan dan saran untuk penelitian selanjutnya.
\end{description}